\documentclass[11pt]{article}

\usepackage[T1]{fontenc}
\usepackage{amsmath,amssymb}
\usepackage{booktabs}
\usepackage{hyperref}

\title{Predicate Formulation Incompatibility: FakeSciencePredicate Conflict Analysis}
\author{Formalization verification}
\date{\today}

\begin{document}
\maketitle

\begin{abstract}
This note documents the incompatibility between the Classical Mechanics
predicate formulation and a conflicting axiomatization supplied by the
FakeSciencePredicate module. Using Lean's module system as a verification
framework, we demonstrate that the predicate approach provides automatic
detection of axiomatizations that cannot coexist. The conflict arises not
from subtle semantic disagreement but from explicit redeclaration, making
the incompatibility decidable and machine-checkable.
\end{abstract}

\section{Setup: Two Competing Predicates}

The foundational Classical Mechanics predicate defines:
\[
  \operatorname{CM}(\mathcal{M}) \iff P_1 \land P_2 \land \cdots \land P_6,
\]
where $P_6$ encodes Newton's second law:
\[
  \sum_{i \in I} f(p, t, i) = m(p)\,\frac{\partial^2 s}{\partial t^2}(p, t).
\]

The FakeSciencePredicate module purports to also define a predicate on the
same namespace (\texttt{ClassicalMechanicsPredicate}), but it does so by
asserting a false axiom:
\begin{verbatim}
  axiom conflict : False
\end{verbatim}

This immediately makes any predicate definition derived from FakeSciencePredicate
inconsistent.

\section{Conflict Detection Mechanism}

\subsection{Redeclaration Conflict}

When Lean attempts to process both modules in sequence, it encounters an error:
\begin{verbatim}
  environment already contains
  'ClassicalMechanicsPredicate.isClassicalMechanicsModel'
  from ClassicalMechanicsPredicate
\end{verbatim}

This indicates that FakeSciencePredicate cannot redefine an identifier that
already exists in the shared namespace.

\subsection{False Axiom Dependency}

Any attempt to derive theorems from FakeSciencePredicate results in a
dependency on the false axiom. Lean's axiom tracking feature reveals:
\begin{verbatim}
  'ClassicalMechanicsPredicate.isClassicalMechanicsModel'
  depends on axioms: [ClassicalMechanicsPredicate.conflict]
\end{verbatim}

This traces back to the source of inconsistency.

\subsection{Proof by Contradiction}

If one attempts to combine both definitions within a single file or import
context, the type checker fails:
\begin{verbatim}
  'ClassicalMechanicsPredicate.isClassicalMechanicsModel'
  has already been declared
\end{verbatim}

Lean refuses to proceed, treating the conflict as a compilation error rather
than a deferred consistency issue.

\section{Why This Matters: The Role of the Predicate}

In the Suppes tradition, a set-theoretic predicate is intended to be a
\textit{definition} that characterizes a class of models. The predicate
succeeds when:
\begin{enumerate}
  \item It is internally consistent (provable theorems do not depend on false
    axioms).
  \item It correctly captures the intended theory (all axioms are present and
    sufficient).
  \item It does not admit vacuous models (systems satisfying the predicate are
    genuinely classical mechanical).
\end{enumerate}

FakeSciencePredicate violates all three criteria:
\begin{enumerate}
  \item It introduces a false axiom, making the namespace itself inconsistent.
  \item It attempts to redefine the predicate, corrupting the original definition.
  \item Any system claimed to satisfy the predicate would satisfy a false
    statement, making the class of models empty (in a consistent fragment).
\end{enumerate}

\section{Computational Verification}

The Lean module system provides automated verification of incompatibility:
\begin{enumerate}
  \item \textbf{Compile-time detection}: Lean's parser and type checker
    immediately flag the conflict, providing the exact line and error.
  \item \textbf{No hidden assumptions}: The incompatibility is not buried in
    subtle proof steps; it is explicit in the module structure.
  \item \textbf{Transitive inconsistency}: Any module that imports both
    \texttt{ClassicalMechanicsPredicate} and \texttt{FakeSciencePredicate}
    becomes inconsistent.
\end{enumerate}

The audit module \texttt{PredicateAudit.lean} formalizes this verification
by attempting to reference both definitions and tracing their axiom
dependencies.

\section{Lessons for Scientific Formalization}

\subsection{Transparency}

The predicate formulation makes the structure of a scientific theory explicit:
the choice of axioms, their logical dependencies, and the resulting model class
are all visible in the Lean source.

\subsection{Decidability}

Incompatibility between rival theories is not a matter of opinion or
interpretation. The Lean type system decides, automatically and unambiguously,
whether two axiomatizations can coexist.

\subsection{Modularity}

The predicate approach allows theories to be developed independently and then
compared. If they are incompatible, that fact is discovered at the moment of
composition, not after extensive proof development.

\section{Conclusion}

The Classical Mechanics predicate formulation, embodied in the explicit
conjunction $P_1 \land P_2 \land \cdots \land P_6$, provides a foundation for
unambiguous scientific theory comparison. When FakeSciencePredicate introduces
a conflicting axiomatization, the Lean system detects the incompatibility
deterministically, supporting the view that scientific predicates are not
merely informal descriptions but formal, machine-checkable structures.

\bibliographystyle{plain}
\bibliography{references}

\end{document}
